% !TEX program = xelatex
% ============================================================
% 全国大学生数学建模竞赛论文骨架（日常训练 / 速写版）
%   · 编译：xelatex main.tex（需 TeX Live，含 ctex 宏包）
%   · 正式提交用官方格式模板（已内置）：
%     00_标准/论文模板/CUMCM_国赛/CUMCMThesis/（见其 README.md，电子版用 withoutpreface）
%   · 使用约定：标注 [AI] 的位置由对应角色卡（prompts/05_章节写手 等）填空，
%     图片放 figs/ 子目录，文件名与 figure_index.csv 中的图号一致。
% ============================================================
\documentclass[UTF8,zihao=-4]{ctexart}
\usepackage[a4paper,margin=2.5cm]{geometry}
\usepackage{amsmath,amssymb,amsfonts}
\usepackage{graphicx}
\usepackage{float}
\usepackage{booktabs}
\usepackage{multirow}
\usepackage{caption}
\usepackage{algorithm}
\usepackage{algpseudocode}
\usepackage[hidelinks]{hyperref}

\captionsetup{font=small,labelsep=quad}
\graphicspath{{figs/}}

\begin{document}

% ---------------- 标题 ----------------
\begin{center}
  {\zihao{3}\bfseries 基于XXX的XXX问题研究}
\end{center}

% ---------------- 摘要 [AI：prompts/06_摘要大师，五要素结构] ----------------
\begin{center}\textbf{摘\quad 要}\end{center}
% 骨架：背景1-2句 → 针对问题一/二/三各一段（模型名+方法+数值结论）→ 亮点1-2句
本文研究……问题。针对……，建立……模型，采用……方法，求解得……（具体数值结论）。
\textbf{关键词}：关键词1；关键词2；关键词3
\newpage

% ---------------- 正文 ----------------
\section{问题重述}
% [AI] 用自己的语言重述题目背景与各问要求，不超过一页
\section{问题分析}
% [AI] 每问的分析思路 + 全文技术路线图（make_flowchart.py 生成，放 figs/）
% \begin{figure}[H]\centering\includegraphics[width=.92\textwidth]{flowchart_route.pdf}
%   \caption{全文技术路线图}\label{fig:route}\end{figure}
\section{模型假设}
% [AI] 6–10 条关键假设，每条附合理性说明（对应 prompts/02 的假设清单）
\begin{enumerate}
  \item 假设一：……；
  \item 假设二：……。
\end{enumerate}

\section{符号说明}
% [AI] 单一事实源：全文新符号先入表再使用；与正文定义严格一致
\begin{table}[H]\centering
  \caption{主要符号说明}\label{tab:notation}
  \begin{tabular}{cll}
    \toprule
    符号 & 含义 & 单位 \\
    \midrule
    $x_{ij}$ & 第 $i$ 供应地到第 $j$ 需求地的运输量 & 吨 \\
    \bottomrule
  \end{tabular}
\end{table}

\section{模型的建立与求解}
\subsection{问题一模型的建立与求解}
\subsubsection{建模思路}
% [AI] 为什么是这个模型（2-3句），呼应问题分析
\subsubsection{模型建立}
% [AI] 公式编号建议按问题前缀：(1-1)(1-2)……
\begin{equation}
  \min Z = \sum_{i}\sum_{j} c_{ij} x_{ij}
  \label{eq:p1-obj}
\end{equation}
其中，$Z$ 为总成本；$c_{ij}$ 为单位运输成本（见符号说明表~\ref{tab:notation}）。
\subsubsection{模型求解}
% [AI] 算法步骤或 algorithm 伪代码；参数取值及依据
\begin{algorithm}[H]
  \caption{求解算法}\label{alg:p1}
  \begin{algorithmic}[1]
    \State 初始化种群；
    \While{未达到最大迭代次数}
      \State 计算适应度；选择、交叉、变异；
    \EndWhile
    \State \Return 最优解；
  \end{algorithmic}
\end{algorithm}
\subsubsection{结果分析}
% [AI] 图先出、文字跟图走；每图“引用+解读”；数字全部来自结果CSV
% \begin{figure}[H]\centering\includegraphics[width=.8\textwidth]{fig_line_convergence.pdf}
%   \caption{算法收敛曲线对比}\label{fig:p1-conv}\end{figure}

\subsection{问题二模型的建立与求解}
% [AI] 结构同问题一；开头补一句与问题一的衔接
\subsection{问题三模型的建立与求解}

\section{敏感性分析}
% [AI] make_chart.py sensitivity 生成；结论句式：“参数±10\%扰动下目标值波动≤±8\%，模型稳健”

\section{模型的评价与推广}
% [AI] 优点（对应题目要求的做得好的点）；缺点必配改进方向；推广到同类问题

\begin{thebibliography}{9}
  \bibitem{ref1} 作者. 文献题目[J]. 期刊, 年份, 卷(期): 页码.
\end{thebibliography}

\newpage
\appendix
\section{附录：求解程序源代码}
% [AI] 代码工程师（prompts/03）交付的可复现脚本，关键步骤注释对应文中公式编号
\end{document}
